
\chapter{Data, Telemetry, and Dataset Scope}
\label{ch:data-telemetry}

Every empirical claim in this dissertation rests on two locked dataset
surfaces: the Germany OPSD archive and the United States EIA-930 feed.  This
chapter documents the signals, their provenance, the telemetry pathways
through which they enter the ORIUS loop, and the fault modes that make
degraded observation an ordinary operating condition rather than an edge case.

% ----------------------------------------------------------------
\section{Germany and United States Dataset Surface}

The ORIUS evaluation spans two geographically and regulatorily distinct
power systems.

\paragraph{Germany (OPSD).}
The Open Power System Data platform aggregates hourly generation, load, and
day-ahead price series from ENTSO-E and the German TSOs.  The locked dataset
covers 2018-10-07 to 2020-09-30 (17\,377 hourly rows), with four primary
targets: \texttt{load\_mw}, \texttt{wind\_mw}, \texttt{solar\_mw}, and
\texttt{price\_eur\_mwh}.  Weather covariates (temperature, irradiance, wind
speed) are sourced from the Open-Meteo archive API and joined on UTC
timestamp.  The full feature matrix contains 94 engineered columns after
calendar encoding and lag construction.

\paragraph{United States (EIA-930).}
The U.S.\ Energy Information Administration Form~930 provides hourly
demand and net generation by source for each balancing authority.  Three
regions are retained:

\begin{table}[ht]
\centering
\caption{Dataset cards for the locked evaluation surfaces.  Source:
\texttt{reports/publication/dataset\_cards.csv}.}
\label{tab:dataset-cards}
\begin{tabular}{llrrrl}
\toprule
\textbf{Key} & \textbf{Region} & \textbf{Rows} & \textbf{Features} &
  \textbf{Coverage} & \textbf{Period} \\
\midrule
DE       & Germany (OPSD)    & 17\,377 & 94  & 100\% & 2018-10 to 2020-09 \\
US\_ERCOT & USA (EIA-930)    & 13\,453 & 62  & 100\% & 2024-07 to 2026-01 \\
US\_MISO  & USA (EIA-930)    & 13\,638 & 114 & 100\% & 2024-07 to 2026-01 \\
US\_PJM   & USA (EIA-930)    & 13\,639 & 62  & 100\% & 2024-07 to 2026-01 \\
\bottomrule
\end{tabular}
\end{table}

Weather covariates for US regions are likewise obtained from Open-Meteo.
Carbon intensity is estimated from the SMARD generation mix (DE) or a
proxy generation-mix model (US).

Complete source Data Cards for the battery witness rows and the broader
ORIUS portability package are provided in Appendix~\ref{app:dataset-cards}.

% ----------------------------------------------------------------
\section{Forecast Targets and Operational Channels}

Each dataset surface exposes three operational forecast targets:

\begin{description}
  \item[\texttt{load\_mw}] Net system load (total demand minus behind-the-meter
    solar), measured at the transmission level.  This is the primary signal for
    arbitrage scheduling.
  \item[\texttt{wind\_mw}] Aggregate onshore and offshore wind generation.
    Wind exhibits high variance and heavy tails, particularly during storm
    fronts, making calibration challenging.
  \item[\texttt{solar\_mw}] Aggregate photovoltaic generation.  Solar follows a
    deterministic diurnal envelope modulated by cloud cover; residual
    uncertainty is sharply bimodal (clear vs.\ overcast).
\end{description}

A fourth target, \texttt{price\_eur\_mwh} (DE only), drives the economic
objective in the dispatch optimiser.  For US regions, a simplified tariff
schedule or locational marginal price proxy is used.

All targets are sampled at hourly resolution and aligned to UTC.  The forecast
horizon is $H = 48$\,h with a stride of 1\,h, yielding a rolling-origin
evaluation with 168 test windows per week.

% ----------------------------------------------------------------
\section{Observed Telemetry Versus Truth Trajectory}

A central design choice in ORIUS is the separation of \emph{truth state}
and \emph{observed state}.  At every control step~$t$:

\begin{itemize}
  \item The \textbf{truth trajectory} $\{z_t\}$ records what actually happened
    at the plant: the true SOC, the true net load, the true renewable
    generation.  This trajectory is never available to the controller in real
    time; it is logged by the CPSBench harness for post-hoc evaluation.
  \item The \textbf{observed trajectory} $\{\tilde{z}_t\}$ is what the
    controller actually sees after telemetry faults have been applied.  It may
    contain dropped readings (NaN), delayed values, sensor spikes, stale
    repetitions, or drifted baselines.
\end{itemize}

The gap $\|z_t - \tilde{z}_t\|$ is the observation error.  All safety metrics
(TSVR, severity) are computed on the truth trajectory, while the controller's
decisions are computed on the observed trajectory.  This dual-trajectory
architecture is the foundation of the CPSBench evaluation protocol
(Chapter~\ref{ch:cpsbench}).

% ----------------------------------------------------------------
\section{Missingness, Corruption, and Signal Degradation}

Telemetry faults fall into six categories, each modelled by a corresponding
CPSBench fault injector:

\begin{table}[ht]
\centering
\caption{Telemetry fault taxonomy and operational causes.}
\label{tab:fault-taxonomy}
\begin{tabular}{lll}
\toprule
\textbf{Fault type} & \textbf{Mechanism} & \textbf{Field cause} \\
\midrule
Dropout       & Packet loss (NaN)       & Cellular dropout, gateway reboot \\
Delay jitter  & Arrival-time shift      & Network congestion, buffering \\
Spike         & Multiplicative outlier  & Sensor transient, ADC overflow \\
Stale sensor  & Frozen reading          & Firmware hang, watchdog failure \\
Drift         & Slow baseline shift     & Sensor ageing, calibration decay \\
Reorder       & Out-of-sequence packets & UDP path diversity, queue inversion \\
\bottomrule
\end{tabular}
\end{table}

These fault types are not mutually exclusive.  The \texttt{drift\_combo}
scenario in CPSBench applies all six simultaneously at configurable severity,
representing a worst-case but operationally plausible field condition.

% ----------------------------------------------------------------
\section{Why Telemetry Faults Are Ordinary, Not Edge Cases}

Field data from commercial BESS installations suggests that telemetry
degradation is the norm rather than the exception:

\begin{itemize}
  \item Cellular uplinks in rural wind-farm sites experience 2--8\% packet
    loss during normal operation, rising to 15--25\% during severe weather.
  \item SCADA gateway reboots occur 1--3 times per month, each causing a
    30--90\,s telemetry blackout.
  \item Sensor drift in current transducers accumulates at 0.1--0.5\% of
    full-scale per year, enough to bias SOC estimates by 2--5\,MWh on a
    100\,MW/400\,MWh system within six months.
  \item Firmware updates on BMS controllers occasionally freeze individual
    cell-group readings for minutes until the watchdog timer resets the
    reporting thread.
\end{itemize}

A safety framework that assumes clean telemetry is therefore unsafe by
construction.  The ORIUS evaluation treats faulted telemetry as the
\emph{default} operating condition and clean telemetry as a special case
(the \texttt{nominal} scenario with zero fault injection).

% ----------------------------------------------------------------
\section{Artifact Contracts and Profile Lock}

Each dataset surface is frozen at a specific commit hash and validated by a
reproducibility contract:

\begin{enumerate}
  \item \textbf{Row-count invariant.}  The ingestion pipeline asserts that
    the row count matches the locked value in
    \texttt{reports/publication/dataset\_cards.csv}.
  \item \textbf{Column schema.}  The feature matrix schema (column names,
    dtypes, nullable flags) is checked against a frozen JSON schema at
    pipeline entry.
  \item \textbf{Target coverage.}  All three forecast targets must have
    100\% non-null coverage in the locked period.  Any gap causes the pipeline
    to abort with a traceable error.
  \item \textbf{Hash verification.}  The SHA-256 hash of the serialised
    Parquet file is compared to the locked value before any model training or
    evaluation proceeds.
\end{enumerate}

These contracts ensure that no drift in upstream data sources can silently
invalidate the empirical results.

% ----------------------------------------------------------------
\section{Scope Boundary}

The dataset scope of this dissertation is intentionally bounded:

\begin{itemize}
  \item \textbf{Included:} hourly bulk-power system signals (load, wind,
    solar, price) for four regions across two continents.
  \item \textbf{Excluded:} sub-hourly (5-min, 1-min) data, distribution-level
    signals, behind-the-meter solar, demand response, and frequency regulation
    signals.
  \item \textbf{Rationale:} the hourly resolution matches the natural
    arbitrage cycle of 2--4\,h duration storage and aligns with the temporal
    granularity of conformal calibration windows.
\end{itemize}

Extending the framework to sub-hourly resolution is discussed in
Chapter~\ref{ch:orius-bench-battery} as part of the benchmark roadmap.

% ----------------------------------------------------------------
\section{Chapter Summary}

This chapter documented the two dataset surfaces (Germany OPSD and US EIA-930)
that underpin every empirical result, the three forecast targets (\texttt{load\_mw},
\texttt{wind\_mw}, \texttt{solar\_mw}), and the six telemetry fault types that
degrade observations between the plant and the controller.  The dual-trajectory
architecture---truth state versus observed state---was introduced as the
foundation for safety evaluation.  Artifact contracts and hash-locked profiles
ensure reproducibility.  Telemetry faults are treated as ordinary operating
conditions, not edge cases, motivating the observation-aware machinery
developed in the following chapters.
